\documentclass[12pt]{article}
\usepackage[margin=1in]{geometry}
\usepackage{titlesec}
\usepackage{hyperref}
\usepackage{booktabs}
\usepackage{array}
\usepackage{parskip}
\usepackage{graphicx}

\title{Design Specification\\
\large QTC Workflow Queue (WFQ) Administrator Application}
\author{Dylan Huang \and Laila Velasquez \and Samuel Acevedo \and Haik Oganesyan \\
\and Ruben Flores \and Slok Patel \and Edwin Rojas \and Jesus Ojeda \\
\and Alexander Diaz \and Zhen Zhao}
\date{Version 1.0 \\ October 17, 2025}

\begin{document}

\maketitle
\newpage
\tableofcontents
\newpage

% ─────────────────────────────────────────────
\section{Overview}
% ─────────────────────────────────────────────

This Design Specification provides a detailed breakdown of every page, screen, component, and
tool within the QTC Workflow Queue (WFQ) Administrator Application. It is intended to serve
as a reference for developers building the UI and for QA testers verifying that each screen
meets the specified requirements.

The application is a web-based, server-side rendered ASP.NET Core MVC application. All pages
are accessible only to users who have been granted administrative access via the
\texttt{?is\_admin=true} URL query string flag.

% ─────────────────────────────────────────────
\section{Global Layout and Branding}
% ─────────────────────────────────────────────

\subsection{Header / Navigation Bar}
\begin{itemize}
    \item \textbf{Location:} Top of every page.
    \item \textbf{Contents:}
    \begin{itemize}
        \item Leidos / QTC Health Services logo (top-left).
        \item Application title: ``Workflow Management System'' (center or adjacent to logo).
        \item ``Administrator Portal'' label (top-right).
    \end{itemize}
    \item \textbf{Branding:} Must adhere to Leidos corporate color palette (purple/white
    primary scheme as shown in mockups).
    \item \textbf{Behavior:} The header is static and does not scroll with the page content.
\end{itemize}

\subsection{Footer}
\begin{itemize}
    \item \textbf{Location:} Bottom of every page.
    \item \textbf{Contents:} Copyright notice and application version number.
\end{itemize}

\subsection{Notification / Toast Messages}
\begin{itemize}
    \item \textbf{Trigger:} Displayed after any create, update, or delete operation.
    \item \textbf{Success:} Green toast — e.g., ``Workflow saved successfully.''
    \item \textbf{Error:} Red toast — e.g., ``A workflow with this name already exists.''
    \item \textbf{Behavior:} Auto-dismisses after 4 seconds or can be manually closed.
\end{itemize}

\subsection{Error Page}
\begin{itemize}
    \item \textbf{Trigger:} Displayed when an unhandled server-side exception occurs.
    \item \textbf{Contents:} A user-friendly message (e.g., ``An unexpected error occurred.
    Please try again or contact your administrator.'') with no technical stack trace visible
    to the user.
\end{itemize}

\subsection{Access Denied Page}
\begin{itemize}
    \item \textbf{Trigger:} Displayed when a user attempts to access the application without
    the \texttt{?is\_admin=true} flag.
    \item \textbf{Contents:} ``Access Denied. You do not have permission to view this page.''
\end{itemize}

% ─────────────────────────────────────────────
\section{Page: Workflow Management Dashboard}
% ─────────────────────────────────────────────

\subsection{Purpose}
The main landing page of the application. Displays a paginated, searchable, and sortable list
of all workflows in the WFQ database.

\subsection{URL}
\texttt{/Workflows/Index} (or \texttt{/})

\subsection{Access Control}
Requires \texttt{[Authorize(Policy = "IsAdmin")]}. Redirects to Access Denied page if not
authorized.

\subsection{UI Components}

\subsubsection{Page Title}
\begin{itemize}
    \item Text: ``Workflow Management''
    \item Subtitle: ``Manage and monitor all workflow processes''
\end{itemize}

\subsubsection{Search Bar}
\begin{itemize}
    \item \textbf{Type:} Text input field.
    \item \textbf{Placeholder:} ``Search workflows...''
    \item \textbf{Behavior:} Filters the workflow list by Name on input change or form
    submission. Sends a GET request with a \texttt{search} query parameter.
    \item \textbf{Max Length:} 100 characters.
\end{itemize}

\subsubsection{Sort Dropdown}
\begin{itemize}
    \item \textbf{Type:} Select/dropdown.
    \item \textbf{Options:} ``Name A-Z'', ``Name Z-A'', ``Newest First'', ``Oldest First''.
    \item \textbf{Behavior:} Re-fetches the workflow list with the selected sort order applied.
\end{itemize}

\subsubsection{Create New Workflow Button}
\begin{itemize}
    \item \textbf{Label:} ``+ Create New Workflow''
    \item \textbf{Style:} Primary button (purple, per Leidos branding).
    \item \textbf{Behavior:} Opens the Create Workflow Modal Dialog (Section 4).
\end{itemize}

\subsubsection{Workflow Table}
\begin{itemize}
    \item \textbf{Columns:}
    \begin{itemize}
        \item \textbf{Name} (clickable link, navigates to Workflow Details page)
        \item \textbf{Description} (truncated if too long)
        \item \textbf{Status} (badge: ``New'' / ``In Progress'' / ``Completed'')
        \item \textbf{Actions} (Edit icon, View Steps icon)
    \end{itemize}
    \item \textbf{Empty State:} If no workflows exist or match the search, display: ``No
    workflows found.''
\end{itemize}

\subsubsection{Pagination Controls}
\begin{itemize}
    \item \textbf{Location:} Below the workflow table.
    \item \textbf{Contents:} ``Items per page'' dropdown (options: 10, 25, 50), page number
    buttons, ``Previous'' and ``Next'' buttons, and a ``Page X of Y'' indicator.
    \item \textbf{Behavior:} Sends a GET request with \texttt{page} and \texttt{pageSize}
    query parameters.
\end{itemize}

% ─────────────────────────────────────────────
\section{Modal: Create / Edit Workflow}
% ─────────────────────────────────────────────

\subsection{Purpose}
A modal dialog used for both creating a new workflow and editing an existing one.

\subsection{Trigger}
\begin{itemize}
    \item \textbf{Create:} Clicking the ``+ Create New Workflow'' button on the Dashboard.
    \item \textbf{Edit:} Clicking the Edit (pencil) icon in the Actions column of the
    workflow table.
\end{itemize}

\subsection{UI Components}

\subsubsection{Modal Header}
\begin{itemize}
    \item \textbf{Create mode:} ``Create New Workflow''
    \item \textbf{Edit mode:} ``Edit Workflow''
    \item Close (X) button in the top-right corner.
\end{itemize}

\subsubsection{Form Fields}
\begin{itemize}
    \item \textbf{Name} (required)
    \begin{itemize}
        \item Type: Text input.
        \item Validation: Required. Must be unique across all workflows. Maximum length TBD
        by Dev Team (see Appendix C of SRS).
        \item Error message: ``Workflow Name is required.'' / ``A workflow with this name
        already exists.''
    \end{itemize}
    \item \textbf{Description} (optional)
    \begin{itemize}
        \item Type: Textarea.
        \item Validation: None required.
    \end{itemize}
\end{itemize}

\subsubsection{Action Buttons}
\begin{itemize}
    \item \textbf{Save Workflow} (primary button) --- Submits the form via POST to
    \texttt{/Workflows/Create} or PUT to \texttt{/Workflows/Edit/\{id\}}.
    \item \textbf{Cancel} (secondary button) --- Closes the modal without saving.
\end{itemize}

\subsubsection{Behavior on Success}
Modal closes. Dashboard workflow list refreshes. Success toast is displayed.

\subsubsection{Behavior on Validation Error}
Modal remains open. Inline validation error messages appear below the relevant fields.

% ─────────────────────────────────────────────
\section{Page: Workflow Details}
% ─────────────────────────────────────────────

\subsection{Purpose}
Displays the full details of a single workflow, including its steps, tasks, and history.

\subsection{URL}
\texttt{/Workflows/Details/\{id\}}

\subsection{UI Components}

\subsubsection{Breadcrumb / Back Navigation}
\begin{itemize}
    \item ``$\leftarrow$ Back to Workflows'' link at the top of the page.
\end{itemize}

\subsubsection{Workflow Header}
\begin{itemize}
    \item Workflow Name (large heading).
    \item Workflow Description (subtitle).
\end{itemize}

\subsubsection{Tab Navigation}
Three tabs: \textbf{Steps}, \textbf{Tasks}, \textbf{History}.

\subsubsection{Steps Tab}
\begin{itemize}
    \item \textbf{Create New Step Button:} ``+ Create New Step'' (opens Create/Edit Step
    Modal).
    \item \textbf{Steps Table:}
    \begin{itemize}
        \item Columns: Step Number, Step Name, Description, Actions (Edit / Delete).
        \item Empty state: ``No steps found. Create your first step to get started.''
    \end{itemize}
\end{itemize}

\subsubsection{Tasks Tab}
\begin{itemize}
    \item \textbf{Tasks Table:}
    \begin{itemize}
        \item Columns: Assigned To, Status, Step, Actions (Complete Task / Reassign).
        \item Empty state: ``No tasks found for this workflow.''
    \end{itemize}
\end{itemize}

\subsubsection{History Tab}
\begin{itemize}
    \item \textbf{History List:}
    \begin{itemize}
        \item Displays task and step history in chronological order.
        \item Each entry shows: timestamp, actor, action taken, and outcome.
        \item Empty state: ``No history available for this workflow.''
    \end{itemize}
\end{itemize}

% ─────────────────────────────────────────────
\section{Modal: Create / Edit Workflow Step}
% ─────────────────────────────────────────────

\subsection{Purpose}
A modal dialog used for both creating a new step within a workflow and editing an existing
step.

\subsection{Trigger}
\begin{itemize}
    \item \textbf{Create:} Clicking the ``+ Create New Step'' button on the Steps tab.
    \item \textbf{Edit:} Clicking the Edit (pencil) icon in the step table Actions column.
\end{itemize}

\subsection{Form Fields}
\begin{itemize}
    \item \textbf{Step Name} (required)
    \begin{itemize}
        \item Type: Text input.
        \item Validation: Required. Must be unique within the workflow.
        \item Error message: ``Step Name is required.'' / ``A step with this name already
        exists in this workflow.''
    \end{itemize}
    \item \textbf{Step Number} (optional)
    \begin{itemize}
        \item Type: Number input.
        \item Validation: Must be unique within the workflow if provided. Auto-assigned if
        left blank.
    \end{itemize}
    \item \textbf{Description} (optional)
    \begin{itemize}
        \item Type: Textarea.
    \end{itemize}
\end{itemize}

\subsection{Action Buttons}
\begin{itemize}
    \item \textbf{Save Step} (primary) --- Submits via POST to
    \texttt{/Workflows/CreateStep} or PUT to \texttt{/Workflows/EditStep}.
    \item \textbf{Cancel} (secondary) --- Closes the modal without saving.
\end{itemize}

% ─────────────────────────────────────────────
\section{Modal: Delete Step Confirmation}
% ─────────────────────────────────────────────

\subsection{Purpose}
A confirmation dialog displayed before a workflow step is deleted, to prevent accidental
data loss.

\subsection{Trigger}
Clicking the Delete (trash) icon in the step table Actions column.

\subsection{UI Components}
\begin{itemize}
    \item \textbf{Warning message:} ``Are you sure you want to delete this step? This action
    is irreversible.''
    \item \textbf{Dependency notice} (if applicable): ``This step has active tasks. They will
    be reassigned to the next step.''
    \item \textbf{Confirm button} (red/danger) --- Submits DELETE to
    \texttt{/Workflows/DeleteStep/\{stepId\}}.
    \item \textbf{Cancel button} (secondary) --- Closes the modal without deleting.
\end{itemize}

% ─────────────────────────────────────────────
\section{Task Management Actions}
% ─────────────────────────────────────────────

\subsection{Reassign Task}
\begin{itemize}
    \item \textbf{Trigger:} Clicking the ``Reassign'' button in the Tasks tab.
    \item \textbf{Behavior:} Opens a small modal or inline form allowing the admin to select
    a new user to assign the task to.
    \item \textbf{Validation:} Cannot reassign a task that is already marked ``Completed.''
    \item \textbf{Endpoint:} PUT \texttt{/Tasks/\{id\}/Assign}
\end{itemize}

\subsection{Complete Task}
\begin{itemize}
    \item \textbf{Trigger:} Clicking the ``Complete Task'' button in the Tasks tab.
    \item \textbf{Behavior:} Marks the task status as ``Completed'' in the database.
    \item \textbf{Validation:} Cannot complete a task that is already ``Completed.''
    \item \textbf{Endpoint:} PUT \texttt{/Tasks/\{id\}/Complete}
\end{itemize}

% ─────────────────────────────────────────────
\section{Form Validation Summary}
% ─────────────────────────────────────────────

All forms in the application apply validation on both the client side (JavaScript) and the
server side (ASP.NET Model validation). The following rules apply globally:

\begin{center}
\begin{tabular}{lll}
\toprule
Field & Rule & Error Message \\
\midrule
Workflow Name    & Required                        & ``Workflow Name is required.'' \\
Workflow Name    & Unique across all workflows     & ``A workflow with this name already exists.'' \\
Workflow Name    & Max length (TBD)                & ``Workflow Name exceeds maximum length.'' \\
Step Name        & Required                        & ``Step Name is required.'' \\
Step Name        & Unique within workflow          & ``A step with this name already exists.'' \\
Step Number      & Unique within workflow (if set) & ``This step number is already in use.'' \\
\bottomrule
\end{tabular}
\end{center}

% ─────────────────────────────────────────────
\section{Accessibility and Compliance}
% ─────────────────────────────────────────────

\begin{itemize}
    \item All form inputs must have associated \texttt{<label>} elements.
    \item All interactive elements must be keyboard-navigable (Tab, Enter, Escape).
    \item Color contrast must meet WCAG 2.1 AA standards.
    \item Modal dialogs must trap focus while open and return focus to the triggering element
    on close.
    \item The application should consider conformance with Section 508 compliance rules.
\end{itemize}

% ─────────────────────────────────────────────
% SNAPSHOT 2 ADDITION
% ─────────────────────────────────────────────
\section{Snapshot 2 Addition: Email Notification Behavior}
\textit{Added by CS3338 Course Team --- Snapshot 2.}

This section describes how the Email Notification System integrates with the existing UI
pages. No new pages are added for this feature --- notifications are triggered silently in
the background as a result of existing user actions.

\subsection{Notification: Task Assigned / Reassigned}
\begin{itemize}
    \item \textbf{Trigger page:} Tasks tab on the Workflow Details page (Section 6).
    \item \textbf{Trigger action:} Admin clicks ``Reassign'' and confirms a new assignee.
    \item \textbf{Behavior:} After the task record is updated in the database, the
    \texttt{INotificationService.SendTaskAssignedAsync()} method is called with the new
    assignee's email address.
    \item \textbf{Email subject:} \texttt{[WFQ] Task Assigned: \{TaskName\}}
    \item \textbf{Email body includes:} Task name, workflow name, step name, and a link to
    the workflow details page.
    \item \textbf{Failure behavior:} If the email fails to send, a warning is logged
    server-side. The UI shows the standard success toast for the reassignment action --- the
    notification failure is not surfaced to the admin.
\end{itemize}

\subsection{Notification: Task Completed}
\begin{itemize}
    \item \textbf{Trigger page:} Tasks tab on the Workflow Details page (Section 6).
    \item \textbf{Trigger action:} Admin clicks ``Complete Task'' and confirms.
    \item \textbf{Behavior:} After the task status is updated to ``Completed'' in the
    database, \texttt{INotificationService.SendTaskCompletedAsync()} is called, notifying
    the workflow administrator.
    \item \textbf{Email subject:} \texttt{[WFQ] Task Completed: \{TaskName\}}
    \item \textbf{Email body includes:} Task name, workflow name, completion timestamp, and
    a link to the workflow details page.
    \item \textbf{Failure behavior:} Same as above --- logged silently, not surfaced to the
    admin.
\end{itemize}

\subsection{Notification: Workflow Step Created}
\begin{itemize}
    \item \textbf{Trigger page:} Steps tab on the Workflow Details page (Section 6), via
    the Create/Edit Step Modal (Section 7).
    \item \textbf{Trigger action:} Admin submits the ``Create New Step'' form successfully.
    \item \textbf{Behavior:} After the new step is persisted to the database,
    \texttt{INotificationService.SendStepCreatedAsync()} is called, notifying the workflow
    administrator.
    \item \textbf{Email subject:} \texttt{[WFQ] New Step Added: \{StepName\} in
    \{WorkflowName\}}
    \item \textbf{Email body includes:} Step name, step number, workflow name, and a link
    to the workflow details page.
    \item \textbf{Failure behavior:} Same as above.
\end{itemize}

\subsection{Notifications Disabled State}
When \texttt{Notifications:Enabled} is set to \texttt{false} in
\texttt{appsettings.json}, the \texttt{INotificationService} methods are called but
immediately return without sending any email. This allows the feature to be safely disabled
in development and testing environments without code changes.

\subsection{No UI Changes Required}
The notification system operates entirely in the background. No new buttons, forms, or pages
are added to the UI as part of this feature. The only visible change is that the existing
success toast messages remain unchanged --- notification delivery is transparent to the
administrator.

% ─────────────────────────────────────────────
% SNAPSHOT 3 ADDITION
% ─────────────────────────────────────────────
\section{Snapshot 3 Addition: Workflow Export Buttons}
\textit{Added by CS3338 Course Team --- Snapshot 3.}

This section describes the UI changes introduced by the Workflow Export / Report Generation
feature. Two new buttons are added to the Workflow Details page. No new pages are created.

\subsection{Export Buttons — Location and Appearance}
\begin{itemize}
    \item \textbf{Location:} Top-right of the Workflow Details page (Section 5), adjacent
    to the existing workflow header. Positioned on the same row as the workflow name and
    description.
    \item \textbf{Buttons:}
    \begin{itemize}
        \item \textbf{``Export PDF''} --- secondary button style (outlined, not filled).
        Icon: document/PDF icon. Color: consistent with Leidos branding.
        \item \textbf{``Export CSV''} --- secondary button style (outlined, not filled).
        Icon: spreadsheet/CSV icon. Color: consistent with Leidos branding.
    \end{itemize}
    \item \textbf{Spacing:} The two export buttons are grouped together with a small gap
    between them, separated from the workflow title by adequate whitespace.
\end{itemize}

\subsection{Export PDF Button}
\begin{itemize}
    \item \textbf{Trigger:} Admin clicks ``Export PDF''.
    \item \textbf{Behavior:} Browser sends a GET request to
    \texttt{/Workflows/ExportPdf?workflowId=\{id\}}. The server generates the PDF and
    returns it as a file download. The browser prompts the user to save or open the file.
    The admin remains on the Workflow Details page.
    \item \textbf{File name format:} \texttt{\{WorkflowName\}\_Report\_\{YYYY-MM-DD\}.pdf}
    \item \textbf{Loading state:} Button shows a spinner and is temporarily disabled while
    the server generates the file, to prevent duplicate requests.
    \item \textbf{Error state:} If generation fails, a red error toast is displayed:
    ``Failed to generate PDF. Please try again.''
\end{itemize}

\subsection{Export CSV Button}
\begin{itemize}
    \item \textbf{Trigger:} Admin clicks ``Export CSV''.
    \item \textbf{Behavior:} Browser sends a GET request to
    \texttt{/Workflows/ExportCsv?workflowId=\{id\}}. The server generates the CSV and
    returns it as a file download. The admin remains on the Workflow Details page.
    \item \textbf{File name format:} \texttt{\{WorkflowName\}\_Tasks\_\{YYYY-MM-DD\}.csv}
    \item \textbf{Loading state:} Same spinner/disabled behavior as the PDF button.
    \item \textbf{Error state:} If generation fails, a red error toast is displayed:
    ``Failed to generate CSV. Please try again.''
\end{itemize}

\subsection{Empty Data Handling}
\begin{itemize}
    \item If the workflow has no steps, the PDF steps table displays: ``No steps have been
    added to this workflow.''
    \item If the workflow has no tasks, the PDF tasks table and the CSV file display:
    ``No tasks available.'' (PDF) or a header row with no data rows (CSV).
    \item Export buttons are always visible and enabled regardless of whether the workflow
    has steps or tasks.
\end{itemize}

\end{document}
